\chapter{Contexte général et état de l'art}

\section*{Introduction}

Ce premier chapitre situe le projet dans son contexte académique et
technologique. Il rappelle d'abord le cadre institutionnel
(l'\Institute), puis présente les fondements théoriques mobilisés :
\emph{cloud computing}, conteneurisation, architecture des applications
web modernes et grands principes de cybersécurité. Cette base
conceptuelle permet de motiver, dans les chapitres suivants, les choix
de conception et d'implémentation retenus.

\section{Cadre institutionnel}

L'\textbf{\Institute} est un établissement d'enseignement supérieur
mauritanien qui forme depuis plusieurs décennies des cadres dans les
domaines de la gestion, de la comptabilité et, plus récemment, du
développement informatique. Le présent travail s'inscrit dans le
parcours de \textbf{Licence en Développement Informatique} et
constitue le \textbf{Projet de Fin d'Études} (PFE) clôturant ce cursus.

L'objectif pédagogique est double : (i)~mettre en pratique de manière
intégrée les compétences acquises durant la formation (programmation,
bases de données, réseaux, systèmes, sécurité) et (ii)~confronter
l'étudiant à un scénario réaliste de cycle complet --- de la
spécification jusqu'au déploiement et à l'audit de sécurité.

\section{Cloud computing}

\subsection{Définition}

Le \textbf{cloud computing} désigne la mise à disposition à la demande,
via Internet, de ressources informatiques (calcul, stockage, réseau,
bases de données, services applicatifs). Le NIST en donne une
définition de référence reposant sur cinq caractéristiques essentielles
\cite{nist-cloud} :

\begin{itemize}
  \item libre-service à la demande (\textit{on-demand self-service});
  \item accès large via le réseau (\textit{broad network access});
  \item mutualisation des ressources (\textit{resource pooling});
  \item élasticité rapide (\textit{rapid elasticity});
  \item facturation à l'usage (\textit{measured service}).
\end{itemize}

\subsection{Modèles de service}

Trois modèles de service principaux sont communément distingués :

\begin{itemize}
  \item \textbf{IaaS} (\textit{Infrastructure as a Service}) : le
        fournisseur expose des ressources virtualisées brutes
        (machines, réseaux, disques). Le client gère l'OS, le \emph{middleware}
        et l'application. \emph{Exemple : AWS EC2.}
  \item \textbf{PaaS} (\textit{Platform as a Service}) : le fournisseur
        gère l'infrastructure et l'environnement d'exécution. Le client
        fournit le code applicatif. \emph{Exemple : AWS Elastic Beanstalk,
        Heroku.}
  \item \textbf{SaaS} (\textit{Software as a Service}) : le client
        consomme directement une application complète accessible via le
        Web. \emph{Exemple : Gmail, Microsoft~365.}
\end{itemize}

Pour ce projet, le modèle retenu est l'\textbf{IaaS}, car il offre le
meilleur compromis entre maîtrise pédagogique des couches techniques
(système, réseau, conteneurs, sécurité) et simplicité de mise en
\oe{}uvre.

\subsection{Le fournisseur AWS}

\textbf{Amazon Web Services} est le leader mondial du \emph{cloud}
public, avec plus d'une trentaine de \emph{régions} (zones
géographiques indépendantes) et une offre couvrant l'ensemble des
modèles de service. Les briques utilisées dans ce projet sont :

\begin{itemize}
  \item \textbf{EC2} (\textit{Elastic Compute Cloud}) : instances de
        machines virtuelles facturées à l'heure ou à la seconde;
  \item \textbf{Elastic IP} : adresse IPv4 publique fixe, associable à
        une instance EC2;
  \item \textbf{Security Group} : pare-feu \emph{stateful} fonctionnant
        au niveau de l'instance;
  \item \textbf{Route\,53} : service DNS managé (zone hébergée et
        enregistrement de domaine).
\end{itemize}

\section{Conteneurisation}

\subsection{Du \emph{bare-metal} aux conteneurs}

L'évolution des modèles d'exécution applicatifs peut se résumer en
trois grandes étapes :

\begin{enumerate}
  \item \textbf{\emph{Bare-metal}} : un serveur physique = une
        application. Coût élevé, ressources sous-utilisées.
  \item \textbf{Virtualisation} : un serveur physique exécute plusieurs
        machines virtuelles, chacune contenant son OS complet. Meilleure
        densité, mais surcharge importante (chaque VM démarre un noyau).
  \item \textbf{Conteneurs} : plusieurs processus isolés partagent le
        \emph{noyau} de l'hôte, en utilisant les \emph{namespaces} et
        les \emph{cgroups} Linux pour cloisonner les ressources et la
        visibilité. Démarrage en quelques centaines de millisecondes,
        empreinte mémoire réduite.
\end{enumerate}

\subsection{Docker et Docker Compose}

\textbf{Docker} est aujourd'hui l'outil de référence pour la création
et l'exécution de conteneurs. Il introduit deux concepts fondamentaux :

\begin{itemize}
  \item l'\textbf{image} : un \emph{instantané} immuable contenant
        l'ensemble des fichiers nécessaires à l'exécution d'une
        application (binaires, bibliothèques, configuration);
  \item le \textbf{conteneur} : une instance d'exécution d'une image,
        avec son propre espace de noms (réseau, processus, fichiers).
\end{itemize}

\textbf{Docker Compose} est un outil d'\emph{orchestration locale} qui
permet, à partir d'un unique fichier déclaratif YAML
(\texttt{docker-compose.yml}), de décrire un ensemble de services, leurs
relations, leurs réseaux et leurs volumes. Une seule commande suffit
ensuite à démarrer ou arrêter l'ensemble de la pile applicative.

\subsection{Bénéfices pour ce projet}

\begin{itemize}
  \item \textbf{Reproductibilité} : le même fichier
        \texttt{docker-compose.yml} fonctionne à l'identique sur le
        poste de développement et sur l'instance EC2.
  \item \textbf{Isolation} : chaque service tourne dans son propre
        conteneur, avec ses dépendances précises.
  \item \textbf{Sécurité} : le réseau interne Docker permet de
        n'exposer que le proxy Nginx, en gardant l'application Flask et
        la base PostgreSQL invisibles depuis Internet.
  \item \textbf{Portabilité} : la migration vers un autre fournisseur
        \emph{cloud} ne nécessite que de relancer la même pile sur une
        nouvelle machine.
\end{itemize}

\section{Architecture des applications web modernes}

\subsection{Le modèle 3-tiers}

Une application web suit classiquement un découpage en trois couches :

\begin{itemize}
  \item \textbf{Présentation} (\emph{front-end}) : navigateur du
        client, qui interprète HTML, CSS et JavaScript;
  \item \textbf{Logique métier} (\emph{back-end}) : serveur applicatif
        qui traite les requêtes, exécute la logique et rend des
        réponses;
  \item \textbf{Données} : base de données qui assure la persistance.
\end{itemize}

\subsection{Le rôle du proxy inverse}

Un \textbf{proxy inverse} est un serveur intermédiaire qui se place
entre les clients et l'application. Il assure plusieurs fonctions
essentielles :

\begin{itemize}
  \item \textbf{Terminaison TLS} : il déchiffre le trafic HTTPS, ce qui
        évite à l'application de gérer elle-même les certificats;
  \item \textbf{Routage} : il dirige les requêtes vers le bon
        \emph{back-end};
  \item \textbf{Mise en cache} et compression;
  \item \textbf{Sécurité périmétrique} : ajout d'en-têtes, limitation
        de débit, filtrage.
\end{itemize}

Pour ce projet, le proxy choisi est \textbf{Nginx}, en raison de sa
légèreté, de sa robustesse et de la richesse de sa documentation.

\section{Cybersécurité applicative}

\subsection{Le Top 10 OWASP}

L'\textbf{Open Web Application Security Project} (OWASP) publie
périodiquement la liste des dix risques applicatifs les plus
critiques. La version 2021, qui sert de référence à ce projet, met
notamment l'accent sur \cite{owasp-top10} :

\begin{enumerate}
  \item \textbf{Broken Access Control} : contournement des contrôles
        d'accès, notamment via la manipulation d'identifiants
        (\textit{IDOR});
  \item \textbf{Cryptographic Failures} : stockage ou transmission de
        données sensibles sans chiffrement adéquat;
  \item \textbf{Injection} : injection SQL, NoSQL, OS command;
  \item \textbf{Insecure Design} : absence de modélisation de la menace
        en amont du développement;
  \item \textbf{Security Misconfiguration} : valeurs par défaut, ports
        inutiles, en-têtes manquants;
  \item \textbf{Vulnerable and Outdated Components} : utilisation de
        bibliothèques non maintenues;
  \item \textbf{Identification and Authentication Failures};
  \item \textbf{Software and Data Integrity Failures};
  \item \textbf{Security Logging and Monitoring Failures};
  \item \textbf{Server-Side Request Forgery (SSRF)}.
\end{enumerate}

\subsection{La défense en profondeur}

La \textbf{défense en profondeur} (\textit{defense-in-depth}) consiste à
empiler plusieurs barrières de sécurité indépendantes, de telle sorte
que la compromission d'une couche ne suffise pas à compromettre le
système entier. Ce principe, hérité du domaine militaire, est
particulièrement pertinent pour les applications hébergées dans le
\emph{cloud}, où les surfaces d'attaque sont multiples (réseau,
système, conteneurs, code, données).

Le chapitre 3 montrera comment ce principe se traduit concrètement dans
WebCyber par six couches de sécurité empilées.

\subsection{TLS et HTTPS}

\textbf{Transport Layer Security} (TLS) est le protocole standard pour
chiffrer les communications sur Internet. \textbf{HTTPS} désigne
simplement HTTP encapsulé dans TLS. Trois propriétés sont garanties :

\begin{itemize}
  \item \textbf{Confidentialité} : le contenu des échanges ne peut être
        lu par un tiers;
  \item \textbf{Intégrité} : le contenu ne peut être modifié sans
        détection;
  \item \textbf{Authenticité} : le client peut vérifier l'identité du
        serveur grâce à un certificat signé par une autorité de
        confiance.
\end{itemize}

\textbf{Let's Encrypt} est une autorité de certification gratuite
fondée en 2014 par l'\emph{Internet Security Research Group}, qui
émet des certificats TLS valides 90 jours, renouvelables
automatiquement via le protocole \textbf{ACME}.

\section{Conclusion du chapitre}

Ce chapitre a présenté les concepts clés mobilisés tout au long du
projet : modèle IaaS et services AWS, conteneurisation par Docker,
architecture web 3-tiers, principes OWASP et défense en profondeur,
ainsi que TLS et Let's Encrypt. Le chapitre suivant s'appuie sur ces
fondations pour formaliser les besoins fonctionnels et non fonctionnels
de l'application WebCyber.
